\documentclass[12pt,a4paper]{article}

\usepackage[utf8]{inputenc}
\usepackage[T1]{fontenc}
\usepackage[margin=2.5cm]{geometry}
\usepackage{hyperref}
\usepackage{microtype}
\usepackage{parskip}
\usepackage{lmodern}

\hypersetup{colorlinks=true, urlcolor=blue!60!black}

\pagestyle{empty}

\begin{document}

\begin{flushright}
Kundai Farai Sachikonye \\
Auf der Lichtung 19, 82178 Puchheim, Germany \\
\href{https://github.com/fullscreen-triangle}{github.com/fullscreen-triangle} \\
\texttt{kundai.sachikonye@bitspark.com} \\
(+49) 1626386344 \\[6pt]
17 May 2026
\end{flushright}

\vspace{1em}

Dr.\ Jia Min Chin \\
Department of Functional Materials and Catalysis \\
Faculty of Chemistry, Universität Wien \\
Vienna, Austria

\vspace{1.5em}

\textbf{Re: University Assistant Predoctoral --- PhD Position in Inorganic \&
Materials Chemistry (Job-ID 5571)}

\vspace{1em}

Dear Dr.\ Chin,

I am applying for the predoctoral position in your group. My background spans
colloidal self-assembly, protein and peptide handling, mass spectrometry
characterisation, and a first-principles framework for predicting the
spectroscopic properties of novel molecular structures without reference
standards --- a combination I believe connects directly to the specific
challenges of field-driven MOF/COF synthesis.

\textbf{Colloidal synthesis and self-assembly.}
During my MSc in Pharmaceutical Biotechnology at HAW Hamburg, I worked
on the encapsulation of proteins inside inverse micelles: selecting the
surfactant system, controlling the water-to-surfactant ratio $w_0$, and
tuning the hydrophilic core environment to stabilise a folded protein in
a non-polar medium. This is supramolecular self-assembly in the same
physical sense as MOF/COF growth --- a cavity with defined chemical
geometry nucleates around a guest molecule or forms through cooperative
non-covalent interactions among building blocks. The specific challenge
of controlling cavity size, polarity, and guest compatibility in inverse
micelles is structurally identical to the pore engineering problem in
framework materials, and the experience of manipulating assembly conditions
--- solvent composition, concentration, temperature, and field effects ---
translates directly.
I have also prepared and characterised lipid nanoparticles and liposomes
(DLS, zeta potential, electron microscopy) and worked on the production
and characterisation of therapeutic lipid colloids at Fraunhofer Institut IGB.
The colloidal assembly dimension of your project is therefore familiar ground.

\textbf{DNA and peptide handling.}
The position lists DNA or peptide handling as a preferred background.
Across several roles I have worked directly with proteins and peptides:
at the University Medical Center Groningen I conducted UPLC-MS/MS shotgun
proteomics, which requires enzymatic protein digestion, peptide fractionation,
and LC separation of complex peptide mixtures;
at Universitätsklinikum Hamburg-Eppendorf I built automated characterisation
pipelines for proteomics and glycomics samples from LC-MS/MS and MALDI-TOF
data, again working with digested proteins and peptide-level identification.
The inverse micelle project at HAW Hamburg involved stabilising folded
proteins inside the micelle core --- handling proteins under conditions
designed to preserve their native structure in a non-aqueous environment.
These roles establish direct hands-on experience with peptides and proteins
as the relevant material class for DNA/peptide-directed MOF/COF synthesis.

\textbf{Spectroscopic characterisation of novel framework structures.}
When a new MOF/COF linker combination is synthesised, no reference IR or
Raman spectrum exists yet in any library. The
\href{https://github.com/fullscreen-triangle/borgia}{Borgia} framework
I have developed addresses this directly: it derives molecular vibrational
modes and harmonic coupling from bounded phase space geometry alone,
requiring no empirical calibration against reference data.
The foundational result is that molecular absorption-emission intervals
create intrinsic frequency ticks; modes sharing rational frequency ratios
$\omega_i/\omega_j = p/q$ (max$(p,q) \leq 10$) form coupled harmonic networks
with closed loops that act as virtual resonant cavities.
This has been validated on six molecules (H$_2$ through C$_6$H$_6$) against
NIST spectroscopic data with self-consistency errors below 1.63\% and zero
adjustable parameters. For a novel MOF linker, the framework predicts which
vibrational modes will harmonically couple, providing theoretical fingerprints
for IR/Raman assignment and identifying which coupling patterns are diagnostic
for specific coordination motifs in the grown framework.
The \href{https://honjo-masamune.vercel.app/tools}{cheminformatics instrument}
built on this framework is available in the browser: zero-parameter
molecular spectroscopy including vibrational modes, docking, 3D structure,
and harmonic network analysis, requiring no installation or account.
A companion paper --- \textit{Harmonic-Scattering Loop Coupling} --- extends this
to a full transfer-matrix framework for coupled molecular resonators,
achieving machine-precision reconstruction across 18/18 benchmarks; and
a \href{https://space-observatory-xi.vercel.app/instruments/resonant-stack/}{Resonant Stack instrument}
demonstrates the physical consequence: Beer--Lambert absorption,
chromatographic separation, and Kirchhoff radiation law are mathematically
identical operations in partition coordinates.
For a novel MOF/COF linker this means the vibrational mode predictions cover
IR absorption, Raman scattering, and thermal emission simultaneously ---
the same first-principles calculation, no re-calibration per spectroscopic method.

\textbf{Field-driven alignment and charge transport anisotropy.}
Your project specifies three field types: electric, magnetic, and flow.
The flow case is addressed directly by the
\href{https://gas-giant.vercel.app/fluids}{Gas Giant} platform I have
built, which derives fluid dynamics from partition network interference ---
Kirchhoff circuit laws map to Navier-Stokes; viscosity emerges as
$\mu = \tau_c \times g$ without phenomenological fitting; Poiseuille flow
profiles validate to 0.9\% mean error.
For flow-directed MOF/COF assembly, the flow field imposes a preferred
partition direction through viscous alignment forces, and the partition
lag $\tau_c$ --- the same parameter that governs viscosity --- determines
the alignment timescale as a function of flow velocity and particle size.
The physical mechanism by which electric and magnetic fields direct growth is
that the field breaks rotational symmetry of the assembly environment,
biasing crystal growth along a preferred axis. In the partition coordinate
framework underlying Borgia, this is equivalent to imposing a preferred
partition direction: the coordination environment along the field axis
develops a deeper, more nested structure than the transverse directions.
Charge transport anisotropy follows because conductivity along the field
axis reflects the higher partition depth of that structural direction.
The composition-inflation theorem ($T(n,d) = d(d{+}1)^{n-1}$) predicts
that the anisotropy ratio scales combinatorially with the number of
field-aligned assembly cycles --- a quantitative, testable connection
between synthesis parameters (field strength, duration) and the resulting
transport tensor that could be validated against the electrochemical
sensing data your group generates.

\textbf{Background.}
I hold an MSc in Pharmaceutical Biotechnology (HAW Hamburg), an MSc in
Computational Biology (Universität Tübingen), and completed doctoral research
in Computational Lipidomics at TUM. I am legally authorised to work in Austria
and am available from June 2026. English is native; German is conversational.

\vspace{1.5em}
Yours sincerely,

\vspace{2.5em}
\textbf{Kundai Farai Sachikonye}

\end{document}
